% ============================================================
%  \section: Multi-Party Conversational AI
%  \section: Research Gaps
%  \section: Positioning of AIutami
%  Parent: \section{State of the Art}
% ============================================================

\section{Multi-Party Conversational AI}
\label{sec:multiparty-ca}

The conversational agents in widespread use today, including Siri, Alexa, ChatGPT, and GPT-4o voice mode, are designed for dyadic interaction: one human speaks to one artificial interlocutor, alternating turns in a strict A-B-A pattern. Real human communication, on the other hand, is mostly group-based. Meetings, classroom discussions, therapy sessions, and brainstorming workshops all involve several participants whose contributions interleave, overlap, and address different subsets of the group.

The move from dyadic to multi-party conversation is a qualitative paradigm shift, not an incremental extension~\cite{gu2022mpc,zheng2022polyadic,traum2004issues}. In a dyadic exchange the information flow is sequential and predictable. In a multi-party setting it becomes a graph: any utterance can come from any participant and can be addressed to any individual, any subset, or the whole group, which gives rise to parallel sub-conversations, shifting coalitions, and complex turn-allocation patterns~\cite{gu2022mpc}. These structural differences propagate to every layer of a conversational AI system, from speaker identification and turn management to response generation and addressee resolution, and they call for dedicated modelling strategies that have no counterpart in dyadic research.

The rest of this section reviews three taxonomies that capture the computational, design, and user-experience dimensions of multi-party conversational AI, and then discusses the technical obstacles that explain why most multi-party research has so far stayed text-based.

\paragraph{Gu et al. (2022): WHO Says WHAT to WHOM.}
Gu et al.~\cite{gu2022mpc} present an IJCAI survey that decomposes multi-party conversation into three interrelated computational problems. \emph{Speaker Modelling} (WHO) is about identifying and predicting the current and next speaker; solutions range from Conditional Random Fields and LSTMs to Transformer-based turn-taking predictors. \emph{Utterance Modelling} (WHAT) addresses what the agent should say, with retrieval-based methods such as TopicBERT alongside generation-based approaches like HeterMPC and, more recently, LLM prompting. \emph{Addressee Modelling} (WHOM) tackles the problem of who an utterance is directed at: addressee recognition, dialogue disentanglement, and the disambiguation of shared-channel messages. Gu et al.\ note that the three problems are usually studied in isolation, but they are deeply interdependent in practice: knowing who speaks conditions the interpretation of what is said, which in turn constrains the inference of whom it addresses. A key limitation of their survey is that it covers only text-based multi-party conversation. Extending it to speech adds further complexity, including ASR errors, overlapping speech, prosodic cues, and real-time latency constraints, none of which the taxonomy addresses.

\paragraph{Houde et al. (2025): designing the agent's participation.}
Houde et al.~\cite{houde2025koala} propose an HCI-oriented framework, developed at IUI~2025, for specifying and controlling the behaviour of an AI agent in group conversations. The framework is organised along two orthogonal axes. The first axis covers \emph{what to control}: WHEN governs the triggers, filters, and frequency of agent contributions (only when addressed, when a self-assessed relevance score exceeds a threshold, or at a fixed cadence); WHAT governs the content, style, and modality of contributions (conservative vs.\ creative ideas, formal vs.\ casual tone); WHERE governs the channel used by the agent (a shared group channel, a threaded reply, or a private message). The second axis covers \emph{how to control it}: SPECIFY describes the mechanism for configuring the agent's behaviour (a UI panel, natural-language instructions, role-based presets); ACCESS describes who is authorised to adjust the configuration (an administrator, any participant, or a democratic consensus); IMPLEMENT describes how the configuration is enforced at the system level (system prompt, external logic, or a hybrid). This design taxonomy is the complement of Gu et al.'s decomposition: where the latter describes the technical sub-problems of multi-party conversation, Houde et al.'s framework describes the design space for governing the agent's participation.

\paragraph{Zheng et al. (2022): challenges and design properties.}
Zheng et al.~\cite{zheng2022polyadic} conduct a CHI literature review of 36 papers on polyadic (multi-party) human-AI interaction, distilling four recurrent challenges. \emph{Inefficient communication} covers topic drift, lack of structure, and difficulty reaching consensus. \emph{Lack of engagement} covers unequal participation and passive members. \emph{Relationship-maintenance barriers} include insufficient emotional awareness and difficulty regulating group affect. The fourth, the \emph{need for building connections}, covers ice-breaking, common-ground establishment, and cross-cultural issues. From these challenges they derive three design properties for a polyadic conversational agent. The agent should be \emph{Visible}: participants must be aware that it is there and what its role is. It should be \emph{Ignorable}: its interventions must be non-intrusive, since too many or too long become counterproductive. And it should be \emph{Accountable}: it must be clear who it is responding to and who controls its behaviour. These three properties impose concrete constraints on system design. In text-based environments, visibility can be achieved through distinct usernames or badges, ignorability through message threading, and accountability through explicit @-mentions. In speech-based environments the same properties take on different meanings: the agent's voice is visible by default since it occupies a shared audio channel, it is hard to ignore since one cannot scroll past a spoken utterance, and it is hard to make accountable without explicit turn-assignment mechanisms.

The amount of speech-based multi-party research is small, and the reason is technical. Addlesee et al.~\cite{addlesee2020asr} provide a systematic evaluation of commercial ASR and speaker-diarisation services in dyadic and multi-party settings, and report performance drops that are often too severe for practical use. Speaker diarisation, the task of attributing each audio segment to the correct speaker, has Diarisation Error Rates between 49\% and 67\% in multi-party settings on commercial platforms, and Google's service fails entirely when more than three speakers are present. Overlapping speech, which is natural whenever speakers compete for the floor or produce backchannels, raises Word Error Rates by about 20\%. Incremental ASR, needed for responsive turn-taking, gives a sixfold increase in WER over batch transcription (33\% versus 5\%). Disfluencies such as filled pauses, self-corrections, and edit terms add further variability across services. These compound difficulties explain why the multi-party conversational AI research surveyed by Gu et al.~\cite{gu2022mpc}, Zheng et al.~\cite{zheng2022polyadic}, and Houde et al.~\cite{houde2025koala} has stayed almost exclusively in the text modality.

% ============================================================
\section{AI Agents and their Roles in Group Conversation}
\label{sec:ai-roles}

The literature on group conversation with AI spans a spectrum of agent roles that mirrors a long-standing distinction in group decision research. At one end of the spectrum is the \emph{facilitator}, defined by Schwarz~\cite{schwarz2017facilitator} as an agent that is substantively neutral, has no authority over the decision itself, and intervenes to help the group improve how it identifies and solves problems. At the other end is the \emph{moderator}, a more directive role that enforces structure, allocates turns, and steers the discussion toward a goal. The systems reviewed below populate this spectrum and show what each role implies in terms of design, turn management, and evaluation.

At the most peer-like end of the spectrum the AI sits alongside human members and contributes ideas of its own. Koala~\cite{houde2025koala} is a Slack-based chatbot deployed in three-person human brainstorming groups at IBM. Houde et al.\ evaluate two variants: in the \emph{reactive} variant Koala contributes only when explicitly addressed via @-mention, while in the \emph{proactive} variant it uses a self-scoring LLM mechanism that estimates the value of a possible contribution and posts when the score exceeds a configurable threshold. The proactive variant generated 73\% of all ideas in the AI-augmented condition but only 33\% of the top-rated ones, and 72\% of participants preferred the reactive variant, describing the proactive agent as ``a pedantic student who wouldn't create space for others''. A follow-up study, Koala~II, added a user-tunable proactivity threshold: no group reverted to fully reactive mode, and the perceived usefulness of the controls was rated 4.46 out of 5. The central finding is that agent proactivity should be neither binary nor static, and that users want to adjust it dynamically during a session.

Closer to the centre of the spectrum, the AI takes a supporting role: it responds when addressed but does not steer the discussion. Two recent robot-based systems illustrate this pattern. The ARI robot~\cite{addlesee2024ari} is deployed in a hospital memory clinic, where it interacts with two human participants (a patient and a companion) via speech. ARI uses a cascaded pipeline with gaze-based addressee detection through an LLM classifier, reaching 85.4\% accuracy compared to 53.4\% with text alone. The system also implements incremental Clarification Requests (iCR) for incomplete utterances, which are common with patients experiencing cognitive decline, and a grounding technique that attributes information to a fictitious persona, reducing hallucination by about 10 percentage points across annotated responses. The Furhat robot~\cite{axelsson2025furhat} is a social robot that engages in open-ended conversation with two human participants. Its pipeline processes voice direction-of-arrival, Azure STT, speaker diarisation, and face recognition, then forwards the result to GPT-3.5 for response generation. Addressee accuracy reaches 92.6\% in parallel settings, drops to 79.3\% in true group conversation, and falls to just 26.8\% with audio alone on concurrent speakers, confirming that audio is not enough to resolve addressee ambiguity. Both robots stay close to dyadic interaction with two human participants, and neither takes a moderating role.

Further along the spectrum, the AI acts as a facilitator: it supports the social process of the group without dominating it, often in emotionally sensitive contexts. The Empathic Co-facilitator~\cite{adikari2022empathic} is a text-based AI co-facilitator for online cancer-support groups (Cancer Chat Canada), analysed across approximately 120\,000 conversational turns. It runs alongside a human therapist, monitors group dynamics, detects emotions using Plutchik's eight-emotion taxonomy via Word2Vec and GloVe embeddings, and models emotion-state transitions through Markov chains. When the group's negativity score exceeds a threshold, the system either generates an empathic response from pre-defined templates or alerts the therapist. The system pre-dates the LLM era, so its language capabilities are limited to template-based responses, and it operates exclusively in asynchronous text. A more recent and directly experimental study by Gao et al.~\cite{gao2025facilitation} compares AI versus human facilitation in peer-support groups for exam anxiety, with 112 participants split between AI-facilitated and human-facilitated discussions. Across all measures, including emotional connection, perceived facilitation effectiveness, perceived utility, and intention to use the system again, human facilitators outperformed the AI. This is one of the few direct experimental comparisons of the two facilitation conditions, and it suggests that AI moderation is not yet a drop-in replacement for human-led group support.

At the most directive end of the spectrum, the AI acts as a full moderator: it allocates turns, enforces structure, and steers the discussion toward a goal. The historical precursor is VIRTUAL HUMAN~\cite{wahlster2023virtual}, a multimodal sports-quiz system from 2006 featuring two human contestants and three virtual agents (one moderator and two domain experts). Dialogue management relied on symbolic rules rather than statistical or neural models. While the system shows that the idea of an AI moderator in multi-party conversation has a two-decade lineage, its pre-LLM architecture makes it very different from contemporary approaches in both flexibility and scalability. A more recent and methodologically closer precedent is the LLM-powered moderator studied by Dubini~\cite{dubini2024multiuser} in a text-based three-person Murder Mystery task adapted from Stasser and Stewart's hidden-profile paradigm. Across 20 sessions recruited on Prolific (10 without moderator, 10 with the LLM moderator), the work reports that the moderator subtly influenced participation dynamics and the focus of the discussion but had limited measurable effect on decision accuracy on the small sample. Dubini framed the system as a baseline rather than a finished design, leaving open the question of how to make AI moderation effective in this kind of task.

The systems above evaluate AI in group conversation on very different tasks and metrics. Koala measures brainstorming productivity and user preference, ARI and Furhat measure addressee accuracy, Empathic and Gao et al.\ measure emotional connection and perceived effectiveness, and Virtual Human is a proof of concept without a controlled evaluation. The choice of task and metric shapes what we can learn about an AI agent in a group. A particularly useful family of tasks for measuring objective group performance, with a long history in social and organisational psychology, are the survival decision-making tasks reviewed in the next section.

